\section{Conclusion}
\label{sec:conclusion}

In this paper, we investigated the generative linguistic patterns of spatial navigation in Large Language Models. By analyzing narratives generated from the \babi dataset, we identified a systematic "spatial grammar" characterized by low vocabulary entropy and high structural rigidity. Our findings show that \gpto and \gptom default to a formulaic, topological mapping style that lacks the descriptive depth and variance of natural human narrative. However, we also demonstrated that this bias is not a fundamental limitation of model capacity, as explicit "human-like" prompting can significantly increase both structural and lexical diversity.

{\bf Key Takeaway}: The default "RLHF tone" extends into the representation of physical space, inducing a robotic spatial grammar that prioritizes efficiency over naturalism.

Future work should focus on developing a benchmark of human-written spatial narratives to serve as a more direct baseline. Additionally, investigating the relationship between these linguistic patterns and the model's internal activations could reveal whether "spatial units" are activated differently during naturalistic vs. formulaic generation. Finally, exploring how these biases affect the performance and user-perception of real-world embodied agents remains a critical next step.
